\scanpage{89}
% archive sheet number: 91

\uline{Théorème} \unsure{Soit|Soient} $k$ un corps alg. clos de car. $\ne 2$, \deleted{Soit $I$ un ens de card 4, d'où $J$ de card. 3, et $V=V(J)$.} \deleted{Considérons}

Considérons le groupoïde des $(X,I\hookrightarrow X)$, où $X$ est une droite projective sur $k$, \deleted{telle} \inserted{$I$ de card 4.} Ce groupoïde est équivalent à celui des couples $(I,\zeta\in Z(J(I))^*)$, en associant à ce couple

\[
X=X(I)\wedge_{\underset{\text{ou }V(I)}{\underbrace{V(J(I))}}}t(\zeta,I)
\qquad\begin{aligned}&\text{où }t(\zeta,I)\text{ est le }V(I)\text{-}\\[-2pt]&\text{torseur }\ I\wedge X(I)_\zeta^{-1}\end{aligned}
\]

$\deleted{\unsure{Z}}I\hookrightarrow X$ est la fibre \inserted{$X_\zeta$} de $X$ en $\zeta$, compte tenu que

\[
\begin{aligned}
X/V(I)&\simeq X(I)/V(I)\simeq\deleted{X}Z(J)\\
X_\zeta&\simeq X(I)_\zeta\wedge t(\zeta,I).
\end{aligned}
\]

\deleted{\unsure{Cor|Si|NB}} Si $I$ est muni d'un pt $i_0$, on a une description de $X$ comme $X(I)\wedge_{V(I)}X(I)_\zeta^{-1}$, c'est isom. à $X(I)$ une fois choisi un sommet de $X(I)$ au dessus de $\zeta$.

De façon plus générale, \deleted{\unsure{pour|pour un|en}} la donnée d'un isom. ``admissible'' $X\simeq X(I)$ équivaut à celle d'un $V$-isom $I\simeq\deleted{\unsure{I|Z}}X(I)_\zeta$.
